% =====================================================================
%  2.4 Atividade com Docker
% =====================================================================
\subsection{Atividade com Docker}\label{sec:docker}

O \textbf{Docker} é uma plataforma de conteinerização: empacota a aplicação com todas as suas
dependências em uma \textbf{imagem} (gerada a partir de um \texttt{Dockerfile}) e a executa como um
\textbf{container}, um processo isolado que compartilha o \textit{kernel} do host. Por isso, um
container é muito mais leve que uma VM (megabytes em vez de gigabytes, segundos em vez de minutos para
iniciar) e roda igual em qualquer máquina \citep{docker2026,muniz2026}.

Na prática, foi criada a VM \texttt{docker-notes} (\texttt{192.168.1.5}), com o Docker instalado pelo
script de provisionamento. O JAR da aplicação, já compilado pelo Ansible no \texttt{app01}, foi copiado
para essa VM (com \cod|vagrant ssh-config| e \texttt{scp} e depois \cod|vagrant upload|), junto com o
\texttt{application.properties}, e a imagem foi construída com o \texttt{Dockerfile} do
Código~\ref{cod:dockerfile}: \texttt{FROM} define a imagem base (um Java enxuto sobre Alpine Linux),
\texttt{RUN} e \texttt{USER} criam e usam um usuário sem privilégios, \texttt{COPY} copia o JAR e a
configuração, \texttt{EXPOSE} documenta a porta e \texttt{ENTRYPOINT} define o comando de
inicialização. A versão do Java da imagem precisa ser compatível com a usada para compilar o JAR
(Java 21); com uma versão menor, o container falha com \texttt{UnsupportedClassVersionError}.

\codigo{docker}{03-docker-lab/Dockerfile}{Dockerfile da aplicação Easy Notes}{dockerfile}

A imagem foi construída e executada com os comandos do Código~\ref{cod:docker-comandos}. A opção
\path{--add-host} permite que o container encontre o banco de dados do \texttt{db01}.

\begin{terminal}[caption={Construção e execução do container},label={cod:docker-comandos}]
docker build -t devops/notes .                             # constrói a imagem
docker run -d --name notes -p 8080:8080 \
  --add-host db01:192.168.1.4 devops/notes                 # inicia o container
docker ps                                                  # container em execução
docker logs -f notes                                       # log da aplicação
\end{terminal}

Também foi visto o
\textit{multi-stage build}, em que um primeiro estágio compila o projeto com Maven e o segundo leva
para a imagem final só o JAR, deixando-a menor e mais segura (arquivo \texttt{Dockerfile.multistage}
no repositório).

\printtela{docker-container}{Imagem construída e container da aplicação em execução}%
  {Terminal na VM docker-notes com \texttt{docker images}, \texttt{docker ps} (container
   \texttt{notes} com status \textit{Up}) e o \texttt{curl} em \texttt{/api/notes} respondendo.}
